% sec-04-investigator-brochure.tex - Investigator Brochure (full IND)
\section{Investigator Brochure}

This section is the investigational new drug (IND) level summary of the
Investigator's Brochure (IB) for the LLM-directed eight-arm robotic
pancreaticoduodenectomy with perioperative daraxonrasib Phase 1 study. It is
provided here as a clearly labeled supporting document under 21 CFR
\S 312.23(a)(5) and \S 312.55, drawing together, in one place and at IND-summary
depth, the physical, chemical, and pharmaceutical properties of the
investigational drug daraxonrasib (RMC-6236); the available nonclinical
pharmacology, pharmacokinetics, and toxicology that underpin the conservative
160 mg first-in-human starting dose; a full description of the investigational
eight-arm collaborative surgical device (56 degrees of freedom, 640 sensor
channels, the per-arm and cumulative force caps, the layered emergency-stop chain,
and the five-vessel no-fly gate); and a structured summary of the data with
explicit guidance for the investigator. The full, standalone Investigator's
Brochure is a separate bound document submitted with this IND and incorporated by
reference; the present section is the IND-level synopsis and the pointer to that
document, consistent with the ReGARDD convention of using a distinct section to
indicate and include a supporting document. The clinical-conduct detail that the
IB summarizes is carried in the study protocol (\S 6 of this IND), the chemistry,
manufacturing, and control information (\S 7), and the pharmacology and
toxicology information (\S 8), and the investigator is directed to those sections
and to the bound IB for the controlling text \cite{phase1guidance,daraxwhipple}.

\subsection*{Supporting Document Status and Scope}

The Investigator's Brochure assembles, for the clinical investigator and the
reviewing institutional review board (IRB), all reference safety and efficacy
information needed to conduct the study and to assess the expectedness of adverse
events. Because the intervention has two inseparable components, an investigational
drug evaluated under 21 CFR part 312 and an investigational device evaluated under
21 CFR part 812 in a single combined IND and IDE submission, the IB has two
coordinated parts: the daraxonrasib drug brochure and a Physical AI device
supplement that functions as the reference safety information for the device and
the on-premises advisory large language model (LLM). An adverse event is judged
unexpected, and therefore potentially reportable on an expedited timeline, when it
is not listed at the observed specificity or severity in the daraxonrasib IB (for
drug events) or in the Physical AI supplement and System Description (for device
events), as specified in \S 8 of this IND and in the protocol assessment section.
The IB is reviewed and reissued at least annually and whenever new safety
information materially changes the risk profile, and currently enrolled
participants are re-consented when such a change occurs. The version of the IB in
effect at the time of each procedure is recorded in the case report form so that
expectedness determinations are traceable to a fixed reference text
\cite{daraxwhipple,cfr312paiuni}.

\subsection*{Physical, Chemical, and Pharmaceutical Properties of Daraxonrasib}

Daraxonrasib (development code RMC-6236) is an orally bioavailable, small-molecule,
RAS(ON) multi-selective (pan-RAS) inhibitor. Its pharmacological mechanism is a
tri-complex mode of action: the molecule binds the active, guanosine triphosphate
(GTP) bound conformation of mutant RAS in a ternary complex with the abundant
intracellular chaperone cyclophilin A, and in doing so it sterically and
allosterically blocks the engagement of downstream effectors, suppressing the
mitogen-activated protein kinase and phosphoinositide 3-kinase signaling that
drives pancreatic ductal adenocarcinoma (PDAC). Because it targets the shared
GTP-bound state rather than a single codon substitution, daraxonrasib is active
across the spectrum of common oncogenic KRAS variants relevant to PDAC, including
the G12D, G12V, and G12R substitutions that define the eligible population of this
study. The agent is administered as an oral tablet, once daily, and is supplied in
the strengths required to compose the three protocol dose levels of 160 mg, 220 mg,
and 300 mg once daily. Tablets are swallowed whole with water and are not crushed,
split, or chewed; they are packaged in child-resistant, light-protective containers
labeled under 21 CFR \S 312.6 with the investigational-use caution statement,
the protocol number, the lot number, and the storage conditions, and are stored at
controlled room temperature in a secured, temperature-monitored,
access-controlled location until dispensed. The detailed chemistry, manufacturing,
and control characterization of the drug substance and drug product, including
appearance, composition, container-closure system, and the stability program that
governs the labeled expiry, is provided in \S 7 of this IND; the present summary
restates only the properties most relevant to the investigator at the point of
care \cite{daraxwhipple,pdac060s2030}.

The clinical and regulatory standing of daraxonrasib supports its use as the drug
component of this first-in-human perioperative study. The agent received FDA
Breakthrough Therapy designation in June 2025 for previously treated metastatic
PDAC harboring a KRAS G12 mutation, and the dose and exposure assumptions of this
protocol are anchored to the pan-KRAS RASolute 302 program, in which 300 mg once
daily is the established recommended Phase 2 dose (RP2D). In the present study the
drug is investigational specifically in the perioperative, curative-intent
surgical context and for its use in combination with the investigational device;
the escalation deliberately begins one and two steps below the established RP2D to
preserve a conservative first-in-human margin in this new clinical setting
\cite{daraxwhipple,natlpaiplatf}.

\subsection*{Nonclinical Pharmacology, Pharmacokinetics, and Toxicology}

The pharmacology of daraxonrasib is summarized in \S 8.1.1 of this IND and rests on
the RAS(ON) tri-complex mechanism described above: selective engagement of the
GTP-bound mutant RAS in complex with cyclophilin A, with downstream suppression of
effector signaling and consequent antiproliferative and pro-apoptotic activity in
KRAS-mutant PDAC models. Drug distribution and pharmacodynamics support a
once-daily oral schedule with sustained target coverage. The pharmacokinetic
behavior most relevant to this protocol is the relationship between the daily dose,
the serum trough concentration, and the 0.5 ng/mL trough threshold that governs the
perioperative pause-and-restart advisory. A 32-iteration deterministic
pharmacokinetic sweep, performed to characterize the perioperative window,
returned a baseline serum daraxonrasib concentration of 0.45 ng/mL at the operative
timepoint across all 32 iterations, that is, below the 0.5 ng/mL trough threshold
in every iteration, consistent with the protocol-mandated pre-operative pause; the
advisory then recommended a postoperative restart on day T+7 in 29 of 32 iterations
(90.6 percent, corresponding to a Grade A fistula with a sub-threshold trough),
on day T+14 in 3 of 32 iterations (9.4 percent, corresponding to a Grade B or
delayed serum rise), and on day T+21 in 0 of 32 iterations. This sweep is
described in the study protocol (\S 6 of this IND) and is summarized in \S 8 and
\S 3.4 \cite{qspmetpancre,fdadigtwinpc,pdacdigtwinp}.

The toxicology integrated summary (\S 8.1.2) states the safety margins that
underpin the conservative DL1 160 mg first-in-human starting dose. The class
toxicity profile of RAS(ON) multi-selective inhibition is dominated by
gastrointestinal effects (diarrhea, nausea) and dermatologic effects (rash), which
are managed symptomatically and through the protocol dose-interruption and
dose-reduction rules; rescue management of these class effects, including
antidiarrheals and aggressive dermatologic care, is specified in the protocol. The
160 mg starting dose sits a full step below the 220 mg intermediate level and two
steps below the 300 mg established RP2D, providing a deliberate margin for the
perioperative surgical setting in which the drug is paused across the operative
window and restarted only on a fistula-grade-keyed and trough-keyed schedule. The
full data tabulation (\S 8.1.3) tabulates the supporting studies; for this
AI-generated submission the supporting nonclinical and in silico evidence base
includes the author's quantitative-systems-pharmacology metastatic-PDAC simulation
with verification, validation, and uncertainty quantification; the artificial
intelligence digital-twin PDAC simulations; the 100,000-patient in silico Phase 3
five-arm PDAC evaluation; and the end-to-end PDAC digital-twin trial work, each
grounded in the references identified in the figure below. The pharmacology and
toxicology integrated-summary structure that organizes \S 8 is reproduced in
Figure 17 \cite{qspmetpancre,fdadigtwinpc,pdacdigtwinp,chatgpt100kp}.

\begin{mermaidfig}
\node[mmgoal] (pt)    at (5.2,0)      {Pharmacology and\\Toxicology (\S 8)};
\node[mmproc] (ph)    at (0,-2.6)     {8.1.1 Pharmacology\\summary and\\conclusions};
\node[mmproc] (tx)    at (5.2,-2.6)   {8.1.2 Toxicology:\\integrated summary};
\node[mmproc] (ft)    at (10.4,-2.6)  {8.1.3 Toxicology:\\full data tabulation};
\node[mmin]   (pd)    at (-2.6,-5.2)  {Mechanism,\\distribution, RAS(ON)\\pharmacodynamics};
\node[mmin]   (dsim)  at (2.0,-5.2)   {Digital-twin and QSP\\simulation evidence};
\node[mmin]   (safem) at (5.2,-5.2)   {Safety margins; DL1\\160 mg conservative\\first-in-human};
\node[mmin]   (sweep) at (10.4,-5.2)  {32-iteration\\perioperative PK sweep\\(0.45 ng/mL baseline)};
\draw[mmedge] (pt) -- (ph);
\draw[mmedge] (pt) -- (tx);
\draw[mmedge] (pt) -- (ft);
\draw[mmedge] (ph) -- (pd);
\draw[mmedge] (ph) to[out=-30,in=120,looseness=0.6] (dsim.north);
\draw[mmedge] (tx) -- (safem);
\draw[mmedge] (ft) -- (sweep);
\end{mermaidfig}
\figcaption{Figure 17. The pharmacology and toxicology integrated-summary structure
for \S 8. The pharmacology summary covers the RAS(ON) mechanism, drug distribution,
and pharmacodynamics; the integrated toxicology summary states the safety margins
underpinning the conservative DL1 160 mg first-in-human starting dose; and the full
data tabulation tabulates the supporting studies, here grounded in the author's
digital-twin and quantitative-systems-pharmacology simulation evidence and the
32-iteration perioperative pharmacokinetic sweep (0.45 ng/mL baseline serum trough
at the operative timepoint across all iterations).}

\subsection*{Investigational Device Description: the Eight-Arm Robotic System}

The investigational device is an on-premises, LLM-directed, eight-arm parallel
coelomic surgical platform that conducts the pancreaticoduodenectomy and generates
the operative and pathology inputs to the perioperative drug-restart advisory. The
on-premises, repository-pinned LLM acts strictly as a second-opinion advisory
oracle held outside the robot vendor kinematic stack and is never an autonomous
controller; full autonomy is prohibited consistent with 21 CFR \S 312.21(e), and
in this Phase 1 study the device operates only at the clinician-controlled
teleoperated and supervised semiautonomous autonomy levels under continuous human
oversight with immediate manual override. The platform carries 56 mechanical
degrees of freedom, composed of eight arms with seven degrees of freedom each
(8 multiplied by 7), and a 640-channel mixed-rate sensor stack of 80 channels per
arm. Force channels sample at 100 kHz and all remaining channels sample at the
10 kHz command rate. Positioning accuracy is 0.05 mm root mean square at a peak tip
velocity of 1,200 mm/s. The cyber-physical safety envelope comprises a per-arm
tip-force cap of no more than 3 N and a cumulative cross-arm tip-force cap of no
more than 18 N enforced at the heartbeat boundary; a 10 kHz heartbeat coordination
bus broadcasting a 64-byte frame on a 100 microsecond per-arm watchdog deadline; an
emergency arm park within 50 microseconds of a missed or out-of-parameter frame; a
cross-arm emergency stop (E-stop) that halts the full platform within 3 ms at a
0.0 N force clamp; and a software-independent hardware E-stop that meets the 21 CFR
\S 312.404 requirement of no more than 500 ms. The device specifications are
summarized in Table 4.1, and the layered safety envelope, including the five-vessel
no-fly gate, is summarized in Table 4.2 \cite{daraxwhipple,onpremwhippl}.

\begin{table}[t]
\caption{Eight-arm LLM-directed robotic Whipple device specifications (Investigator's Brochure device supplement; 21 CFR \S 312.23(g)(iii) hardware fields).}
\label{tab:sec04-device}
\centering
\begin{tabularx}{\textwidth}{L{4.4cm} L{4.3cm} Y}
\toprule
\textbf{Specification} & \textbf{Value} & \textbf{Notes} \\
\midrule
Architecture & Eight-arm parallel coelomic platform & On-premises LLM advisory oracle outside vendor kinematic stack \\
Degrees of freedom & 56 total (8 arms x 7 DOF) & Cooperating arms, phase-keyed tool assignment P1 to P8 \\
Sensor channels & 640 total (80 per arm) & Force at 100 kHz; all other channels at 10 kHz command rate \\
Positioning accuracy & 0.05 mm RMS at 1,200 mm/s & Verified in the pre-procedure safety matrix each case \\
Per-arm tip-force cap & 3 N maximum & Enforced continuously per arm \\
Cumulative force cap & 18 N maximum & Enforced at the heartbeat boundary across all arms \\
Heartbeat bus & 10 kHz, 64-byte frame & 100 microsecond per-arm watchdog deadline \\
Emergency arm park & 50 microseconds maximum & On a missed or out-of-parameter frame \\
Cross-arm E-stop & 3 ms maximum, 0.0 N clamp & Software escalation halting the full platform \\
Hardware E-stop & 500 ms maximum & Software-independent; meets 21 CFR \S 312.404 \\
Autonomy level (Phase 1) & Teleoperated / supervised & Continuous human oversight; full autonomy prohibited (\S 312.21(e)) \\
\bottomrule
\end{tabularx}
\figcaption{Table 4.1. Device specifications carried in the Investigator's Brochure
device supplement and the Physical AI System Description; the controlling
hardware, software, and safety-system text is in \S 6 and \S 7 of this IND and in
the bound IB.}
\end{table}

The eight cooperating arms are coordinated by the 10 kHz heartbeat bus and carry
phase-keyed tool assignments across the eight operative phases. Phases P1 through
P4 comprise the dissection and specimen removal: Arms 1 and 2 perform the active
dissection under the vessel control of Arm 3 (bipolar plus retractor) and the
near-infrared indocyanine-green imaging of Arm 4, with Arm 7 providing suction and
irrigation throughout. The specimen is removed at the end of P4 and reconstruction
begins. The three reconstructive anastomoses are then constructed in sequence under
closed-loop ring-tension control by the Arm 5 anastomosis probe, which issues a
soft warning whenever measured ring tension falls outside its controller band: the
duct-to-mucosa pancreaticojejunostomy (PJ) in P5 with a band of 0.30 to 0.60 N, the
end-to-side hepaticojejunostomy (HJ) in P6 with a band of 0.20 to 0.50 N, and the
antecolic gastrojejunostomy (GJ) in P7 with a band of 0.40 to 0.80 N; P8 covers
final imaging, hemostasis, and drain placement. The pancreaticojejunostomy is the
dominant determinant of postoperative pancreatic-fistula risk and is therefore the
anastomosis whose ISGPS fistula grade supplies the input to the perioperative
drug-restart advisory \cite{daraxwhipple,paisitedocpk}.

\subsection*{Safety Envelope: Force Caps, E-Stop Chain, and the No-Fly Gate}

The device safety envelope is enforced before and during every procedure rather
than reconstructed afterward. The vascular safety-zone gate ticks at 10 kHz and
evaluates the instrument-tip proximity to each of five named vessels, the superior
mesenteric vein (SMV), the portal vein (PV), the hepatic artery (HA), the celiac
axis (CA), and the superior mesenteric artery (SMA), against three concentric
thresholds: a soft-warning zone at 3.0 mm that raises an advisory but does not block
the command, a no-fly zone that blocks the command (1.0 mm at the SMV and PV; 1.5 mm
at the HA, CA, and SMA), and a hard-stop zone at 5.0 mm that forces a cross-arm
E-stop within 3 ms. Across a representative 100-tick verdict sample the gate
returned clear on 83 ticks, soft warning on 6, no-fly on 6, and hard-stop on 5,
exercising all five vessels across all eight operative phases. The halt chain that
backs the gate is layered: an on-time heartbeat frame continues the command state,
whereas a missed or out-of-parameter frame parks the affected arm within 50
microseconds and escalates to the cross-arm E-stop, and the software-independent
hardware E-stop backs the software chain at the 500 ms outer bound. The safety
envelope is summarized in Table 4.2 \cite{daraxwhipple,clintrustpai}.

\begin{table}[t]
\caption{Device safety envelope: force caps, layered emergency-stop chain, and the five-vessel no-fly gate.}
\label{tab:sec04-safety}
\centering
\begin{tabularx}{\textwidth}{L{3.6cm} L{4.6cm} Y}
\toprule
\textbf{Safety element} & \textbf{Threshold / action} & \textbf{Basis and verdict sample} \\
\midrule
Per-arm force cap & 3 N maximum per arm & Continuous; force channels at 100 kHz \\
Cumulative force cap & 18 N maximum across arms & Enforced at the heartbeat boundary \\
Heartbeat and watchdog & 10 kHz bus, 100 microsecond deadline & On-time frame continues command state \\
Arm park & 50 microseconds maximum & Triggered by a missed or out-of-parameter frame \\
Cross-arm E-stop & 3 ms maximum, 0.0 N clamp & Software escalation; halts full platform \\
Hardware E-stop & 500 ms maximum & Independent of software; 21 CFR \S 312.404 \\
Soft-warning zone & 3.0 mm, all five vessels & Advisory only; 6 of 100 ticks \\
No-fly zone (block) & 1.0 mm SMV/PV; 1.5 mm HA/CA/SMA & Command blocked; 6 of 100 ticks \\
Hard-stop zone & 5.0 mm, all five vessels & Forces E-stop; 5 of 100 ticks \\
Clear verdict & Outside all zones & Command proceeds; 83 of 100 ticks \\
\bottomrule
\end{tabularx}
\figcaption{Table 4.2. The cyber-physical safety envelope and the five-vessel
no-fly gate (SMV, PV, hepatic artery, celiac axis, SMA), sampled at 10 kHz; the
representative 100-tick verdict sample partitions into 83 clear, 6 soft warning,
6 no-fly, and 5 hard-stop ticks. The controlling text is in the protocol assessment
section (\S 6 and \S 8 of this IND).}
\end{table}

Every advisory command, every safety-gate verdict, every force and trajectory
sample at the recorded rates, and every E-stop event is written to a 21 CFR part 11
hash-chained, append-only audit trail, so that compliance with the cyber-physical
envelope is verifiable for each case and any deviation is time-stamped and
detectable. The audit trail also satisfies the preservation requirement from 24
hours before to 72 hours after any reportable Physical AI adverse event. Before any
patient-facing robotic activity, the platform must clear the Phase 0 simulation
validation gate (a Unified Safety Level of at least 7.0; at least 1,000 simulations
across at least two simulation frameworks; a sim-to-real trajectory gap below 2 mm;
and a tip-force gap below 0.5 N), and the gate is re-cleared after any change. The
on-premises LLM is hash-pinned to a deterministic seed and repository commit and is
governed by a verification-before-generation gate suite, aligned with H.R. 9510,
the Verification Before Generation in Physical AI Oncology Trials Act of 2026
\cite{congress9510,clintrustpai,paiautospons}.

\subsection*{Summary of Data and Guidance for the Investigator}

The investigator should treat the daraxonrasib drug brochure and the Physical AI
device supplement as the two controlling reference texts for expectedness and risk.
For the drug, the established 300 mg once-daily RP2D and the class toxicity profile
(gastrointestinal and dermatologic effects) define the expected events; the study
deliberately starts at 160 mg once daily to preserve a conservative first-in-human
margin in the perioperative setting, and the 28-day dose-limiting toxicity window,
the 3+3 escalation, and the dose-interruption and dose-reduction rules govern
management. For the device, the force caps, the layered E-stop chain, and the
five-vessel no-fly gate define the expected safety behavior; any robotic
malfunction, AI prediction error, sensor failure, communication loss, unauthorized
access, digital-twin divergence, or event requiring activation of the E-stop is a
Physical AI adverse event under 21 CFR \S 312.32(f) and is assessed against this
supplement for expectedness. The perioperative pause-and-restart advisory is
LLM-bound and human-confirmed and never autonomous; the investigator confirms the
restart day keyed to the pancreaticojejunostomy ISGPS fistula grade and the serum
trough relative to 0.5 ng/mL, with a Grade C fistula mapping to T+21 or a continued
hold and to the drug stopping rule. All adverse events are graded by CTCAE version
5.0 and reported on the regulatory timelines (no later than 7 calendar days for an
unexpected fatal or life-threatening suspected adverse reaction and no later than
15 calendar days for any other serious and unexpected suspected adverse reaction),
with the six Physical AI reporting triggers of \S 312.32(g) layered on top
\cite{daraxwhipple,cfr312paiuni,ichpaistduni}.

The investigator is reminded that the controlling clinical-conduct text is the
study protocol and that the present section is an IND-level summary; the full
standalone Investigator's Brochure is a separate bound document submitted with this
IND and is the authoritative reference safety information for both the drug and the
device. The investigator should hold a current copy of the bound IB, record its
version in the case report form for each procedure, and consult \S 6, \S 7, and
\S 8 of this IND for the drug-and-device intervention detail, the chemistry,
manufacturing, and control information, and the pharmacology and toxicology
information that the IB summarizes \cite{phase1guidance,daraxwhipple,oncotrialllm}.
